
% ============================================================
% Maxwell--SST Kinetic Closure and Internal-Mode Thermodynamic Gate
% Target branch: SST_CANON-v0.8.35 research track
% ============================================================
\section{Maxwell--SST Kinetic Closure and Internal-Mode Thermodynamic Gate}
    \label{sec:rt_maxwell_sst_kinetic_gate}

\subsection{Status, scope, and source discipline}

\textbf{[ORTHODOX INPUT].}
Maxwell's 1860 kinetic-theory programme derives statistical velocity laws,
collision rates, mean paths, pressure and transport from microscopic moving
bodies, and then extends the calculation to non-spherical bodies whose
collisions exchange translational and rotational energy
\cite{Maxwell1860GasesI,Maxwell1860GasesIIIII}.

\textbf{[DERIVED HISTORICAL LESSON].}
The crucial negative result for the present SST programme is methodological:
when an assumed internal degree of freedom is dynamically coupled and reaches
equilibrium, it must be included in the thermodynamic bookkeeping.  Maxwell
used the resulting specific-heat mismatch as evidence that the mechanical model
was incomplete rather than suppressing the unwanted degree of freedom by
assumption.

\textbf{[SPECULATIVE / SST RESEARCH PROGRAMME].}
The Maxwell particles are \emph{not} identified with elements of the SST
substrate.  The SST substrate remains a continuous incompressible, inviscid
medium.  The proposed kinetic layer applies only to an ensemble of coherent
swirl-string objects or excitations embedded in that medium.

The research question is
\begin{quote}
For every SST knot $K$: which translational, orientational, Kelvin, twist,
writhe, and core modes are actually coupled by a physical interaction; what are
their excitation thresholds or gaps; and what thermodynamic and spectroscopic
contribution follows?
\end{quote}

\subsection{Priority I: internal-mode thermodynamic gate}
    \label{subsec:rt_maxwell_mode_gate}

\subsubsection{Reference state and mode coordinates}

Let $\mathcal Q_K^{(0)}$ denote a preregistered relaxed finite-core state of
knot class $K$.  A local perturbation is expanded in generalized coordinates
$q_{K,a}$,
\begin{equation}
    \mathcal Q_K
    =
    \mathcal Q_K^{(0)}
    +
    \sum_a q_{K,a}\,\mathbf e_{K,a}
    +
    O(q^2).
    \label{eq:rt_maxwell_mode_expansion}
\end{equation}

\textbf{[BRIDGE / RESEARCH DEFINITION].}
For a declared resolved energy functional $E_K[\mathcal Q_K]$, define the
quadratic stiffness and inertia matrices by
\begin{align}
    H^{(K)}_{ab}
    &:={}
    \left.
    \frac{\partial^2 E_K}{\partial q_{K,a}\partial q_{K,b}}
    \right|_{\mathcal Q_K^{(0)}},
    \label{eq:rt_maxwell_hessian}\\
    M^{(K)}_{ab}
    &:={}
    \text{generalized inertia matrix of the same declared model}.
    \label{eq:rt_maxwell_inertia_matrix}
\end{align}
The linearized mode problem is then
\begin{equation}
    H^{(K)}\mathbf e_{K,n}
    =
    \omega_{K,n}^{2}
    M^{(K)}\mathbf e_{K,n}.
    \label{eq:rt_maxwell_generalized_eigenproblem}
\end{equation}
Equation~\eqref{eq:rt_maxwell_generalized_eigenproblem} is a generic small-
oscillation construction; it becomes an SST prediction only after $E_K$ and
$M^{(K)}$ are derived from the declared finite-core dynamics.

\subsubsection{Do not confuse frequency with an energy gap}

\textbf{[ORTHODOX INPUT / CRITICAL GUARD].}
For a classical quadratic mode,
\begin{equation}
    E_{K,n}(A)
    =
    \frac12 \kappa_{K,n}A^2 + O(A^3),
    \label{eq:rt_maxwell_classical_mode_energy}
\end{equation}
so arbitrarily small amplitudes imply arbitrarily small excitation energies.
Therefore
\begin{equation}
    \omega_{K,n}>0
    \not\Rightarrow
    \Delta_{K,n}>0 .
    \label{eq:rt_maxwell_frequency_gap_noimplication}
\end{equation}
This is a direct guard against treating torsional rigidity, curvature stiffness,
or a nonzero Kelvin frequency as a discrete thermodynamic gap.

\textbf{[BRIDGE / DEFINITION].}
A true SST excitation gap must instead be defined from the admissible state
space,
\begin{equation}
    \Delta_{K,a}
    :=
    \inf_{\mathcal Q\in\mathcal A_{K,a}^{\rm exc}}
    \left[
        E_K[\mathcal Q]-E_K[\mathcal Q_K^{(0)}]
    \right],
    \qquad
    \Delta_{K,a}>0,
    \label{eq:rt_maxwell_true_gap_definition}
\end{equation}
where $\mathcal A_{K,a}^{\rm exc}$ must exclude the ground-state neighbourhood
for a physically derived reason---for example discrete occupancy, a topological
barrier, a nonlinear activation threshold, or another explicitly stated
selection rule.  No numerical value of $\Delta_{K,a}$ is canonized here.

\subsubsection{Mode taxonomy and double-counting guards}

The following classification is the required audit basis.

\begin{center}
\begin{tabularx}{\linewidth}{p{0.16\linewidth}p{0.22\linewidth}p{0.22\linewidth}X}
\hline
Mode family & Minimal coordinate & Physical coupling channel & Canon guard \\
\hline
translation & $\mathbf X_K,\mathbf P_K$ & net impulse / force & collective motion; not an internal gap \\
orientation & $R_K\in SO(3)$ & net torque / anisotropic encounter & the static orientation coordinate is a symmetry zero-frequency direction in an isotropic unbounded background; rotational kinetic energy and angular momentum, when retained, form a separate collective rotational sector \\
Kelvin/shape & normal centerline perturbations $q^{\rm K}_{K,n}$ & projected hydrodynamic force & genuine internal shape sector; spectrum must converge with spatial resolution \\
twist & framed-tube variables $q^{\rm Tw}_{K,n}$ & torque / frame stress & undefined in a centerline-only model; requires a material frame and finite-core closure \\
writhe & $Wr[\gamma_K]$ & induced indirectly through shape change & \emph{not} an independent degree of freedom by default; treat as an observable projection of shape modes to avoid double counting \\
core & radial/elliptic/profile variables $q^{\rm core}_{K,n}$ & pressure / strain / local core forcing & absent from singular-filament models; requires a resolved core \\
\hline
\end{tabularx}
\end{center}

\textbf{[ORTHODOX GEOMETRIC INPUT].}
The writhe $Wr$ is a functional of the centerline geometry, while framed-tube
self-linking separates into writhe and twist under the declared ribbon
assumptions.  Consequently a shape eigenmode that changes $Wr$ must not be
counted again as an independent ``writhe oscillator'' unless a reduced model
proves that such a coordinate is dynamically independent
\cite{Moffatt1969,Calugareanu1961,White1969}.

\subsubsection{Interaction-to-mode coupling matrix}

Consider a physical interaction $\mathcal I$---for example a two-knot encounter,
an imposed transverse packet, or a background velocity gradient.  Define the
linear coupling amplitude
\begin{equation}
    \mathcal G^{(\mathcal I)}_{K,a}
    :=
    \mathbf e_{K,a}^{\mathsf T}
    \mathbf Q_{\mathcal I},
    \label{eq:rt_maxwell_mode_coupling_projection}
\end{equation}
where $\mathbf Q_{\mathcal I}$ is the generalized-force covector in the same
coordinate basis as $\mathbf q_K$.  This pairing has units of generalized work
per unit mode amplitude; only its vanishing/non-vanishing and a separately
normalized dimensionless coupling are used as selection diagnostics.

For two interacting knots $K$ and $L$, the bilinear transfer matrix
\begin{equation}
    G^{KL}_{ab}
    :=
    \left.
    \frac{\partial^2 V_{KL}}
    {\partial q_{K,a}\partial q_{L,b}}
    \right|_{0}.
    \label{eq:rt_maxwell_pair_coupling_matrix}
\end{equation}

\textbf{[BRIDGE / RESEARCH DEFINITION].}
A mode is called \emph{directly coupled} under $\mathcal I$ only if the
corresponding projection is nonzero after symmetry reduction and numerical
convergence:
\begin{equation}
    \mathcal G^{(\mathcal I)}_{K,a}\neq0 .
    \label{eq:rt_maxwell_direct_coupling_gate}
\end{equation}
A merely kinematically available eigenmode with zero projected coupling is not
thermodynamically activated by that interaction.

\subsubsection{Excitation and equilibration timescale}

\textbf{[ORTHODOX KINETIC FORM / SST BRIDGE].}
If an inelastic encounter cross section $\sigma^{KL}_{a,\rm inel}$ can be
operationally defined for mode $a$, the dilute-ensemble null model gives
\begin{equation}
    \tau^{-1}_{K,a}
    =
    \sum_L n_L
    \left\langle
        v_{\rm rel}\,
        \sigma^{KL}_{a,\rm inel}
    \right\rangle .
    \label{eq:rt_maxwell_mode_rate}
\end{equation}
The dimensions are
$[n_L][v_{\rm rel}][\sigma]={\rm m^{-3}m\,s^{-1}m^2}={\rm s^{-1}}$.
Equation~\eqref{eq:rt_maxwell_mode_rate} is a null model, not a theorem for
long-range vortex interactions.

A mode contributes to an observed equilibrium only when the three gates
\begin{equation}
    \boxed{
    \mathcal G_{K,a}\neq0,
    \qquad
    E_{\rm drive}\gtrsim\Delta_{K,a},
    \qquad
    \tau_{K,a}\lesssim t_{\rm obs}
    }
    \label{eq:rt_maxwell_three_gate_condition}
\end{equation}
are simultaneously satisfied.

\subsubsection{Thermodynamic contribution}

\textbf{[ORTHODOX STATISTICAL-MECHANICS BENCHMARK].}
If the accessible SST internal states form a discrete spectrum
$E_{K,s}=E_{K,0}+\Delta E_{K,s}$ and are genuinely in thermal equilibrium,
the internal canonical partition function is
\begin{equation}
    Z_{K,\rm int}(\beta)
    =
    \sum_s g_{K,s}
    e^{-\beta\Delta E_{K,s}},
    \qquad
    \beta=(k_BT)^{-1},
    \label{eq:rt_maxwell_internal_partition}
\end{equation}
with
\begin{align}
    U_{K,\rm int}
    &=
    -\frac{\partial}{\partial\beta}
    \ln Z_{K,\rm int},
    \label{eq:rt_maxwell_internal_energy}\\
    C_{V,K}^{\rm int}
    &=
    k_B\beta^2
    \frac{\partial^2}{\partial\beta^2}
    \ln Z_{K,\rm int}.
    \label{eq:rt_maxwell_internal_heat_capacity}
\end{align}
These are orthodox benchmark equations, not a claim that the SST substrate has
a thermodynamic temperature \cite{PathriaBeale2011}.

For a single isolated gapped level with degeneracy $g$ and
$\beta\Delta\gg1$,
\begin{equation}
    C_V^{\rm int}
    =
    O\!\left(
        k_B(\beta\Delta)^2e^{-\beta\Delta}
    \right),
    \label{eq:rt_maxwell_lowT_gap_suppression}
\end{equation}
whereas an ungapped classical quadratic degree of freedom that equilibrates
contributes an order-$k_B$ term by equipartition.  This is the Maxwell-style
thermodynamic mode-count gate.

\subsubsection{Spectroscopic contribution}

\textbf{[BRIDGE / RESEARCH BOUND].}
For an external atomic or resonant state $n$, let $\lambda_{n;K,s}$ denote a
dimensionless effective coupling weight connecting that observable to internal
SST state $s$.  A conservative bookkeeping bound is
\begin{equation}
    |\delta E_n|
    \le
    \sum_{K,s}
    |\lambda_{n;K,s}|\,
    p_{K,s}\,
    \Delta E_{K,s},
    \label{eq:rt_maxwell_spectroscopic_bound}
\end{equation}
with $p_{K,s}$ the physically derived occupation probability.  The model must
also predict line broadening, sidebands, or inelastic thresholds whenever the
same coupling opens real transitions.  A fitted energy shift without the
corresponding transition channels is not sufficient closure.

\subsubsection{Maxwell thermodynamic falsifier}

For a preregistered knot candidate $K$, define the experimentally relevant
active set
\begin{equation}
    \mathcal A_K(T,t_{\rm obs})
    :=
    \left\{
        a:
        \mathcal G_{K,a}\neq0,
        \ \Delta_{K,a}\lesssim k_BT,
        \ \tau_{K,a}\lesssim t_{\rm obs}
    \right\}.
    \label{eq:rt_maxwell_active_mode_set}
\end{equation}

\textbf{[DERIVED / CONDITIONAL FALSIFIER].}
If the calculated set $\mathcal A_K$ contains modes whose orthodox
thermodynamic or spectroscopic signatures exceed the empirical upper bounds,
then at least one of the following must fail: the knot assignment, the mode
spectrum, the coupling model, or the assumed equilibration mechanism.  The
constraint may not be repaired by deleting an already-coupled mode after the
comparison target is exposed.

\subsection{Priority II: kinetic closure for ensembles of swirl strings}
    \label{subsec:rt_maxwell_knot_kinetics}

\subsubsection{Knot phase space}

For each topological class $K$, introduce a distribution
\begin{equation}
    f_K
    =
    f_K(
        \mathbf x,
        \mathbf P,
        R,
        \mathbf q,
        \boldsymbol\pi,
        t
    ),
    \label{eq:rt_maxwell_knot_distribution}
\end{equation}
where $R\in SO(3)$ is orientation and $(\mathbf q,\boldsymbol\pi)$ denotes the
resolved internal-mode coordinates and conjugate generalized momenta.  Chirality
or circulation branch labels are discrete indices attached to $K$, not hidden
continuous fit parameters.

\textbf{[ANSATZ / RESEARCH KINETIC CLOSURE].}
A minimal evolution equation is
\begin{equation}
    \frac{\partial f_K}{\partial t}
    +
    \{f_K,H_K\}
    =
    \sum_L
    \left(
        \mathcal C_{KL}^{\rm ideal}
        +
        \mathcal C_{KL}^{\rm rec}
    \right),
    \label{eq:rt_maxwell_sst_boltzmann}
\end{equation}
where the left-hand side is the Hamiltonian/Liouville transport of the chosen
reduced model.

The collision/encounter operator is split because ideal Euler evolution and
reconnection are not the same dynamical sector:
\begin{align}
    \mathcal C_{KL}^{\rm ideal}:
    &\quad
    K,L,\Gamma_K,\Gamma_L
    \ \text{preserved},
    \label{eq:rt_maxwell_ideal_collision_sector}\\
    \mathcal C_{KL}^{\rm rec}:
    &\quad
    \text{topology-changing events allowed only under an explicitly non-ideal
    closure}.
    \label{eq:rt_maxwell_reconnection_collision_sector}
\end{align}
The ideal operator may exchange translation, orientation and internal-mode
energy while preserving the declared invariants of the underlying model.

\subsubsection{Maxwellian distribution as null hypothesis only}

\textbf{[ORTHODOX NULL MODEL / NOT CANONIZED FOR SST].}
For a dilute ensemble of weakly interacting objects with sufficient mixing, the
translational equilibrium null model is Maxwellian,
\begin{equation}
    f_{K,\rm tr}^{(0)}
    \propto
    \exp\!\left[
        -\frac{
        (\mathbf P-M_K\mathbf U)^2
        }{2M_Kk_BT}
    \right].
    \label{eq:rt_maxwell_translational_null}
\end{equation}
Maxwell's derivation motivates this benchmark
\cite{Maxwell1860GasesI}.  SST vortex knots, however, can interact through
long-range induced velocity fields and orientation-dependent couplings, so
Eq.~\eqref{eq:rt_maxwell_translational_null} is a falsifiable null model rather
than an assumed SST equilibrium law.

\subsection{Priority III: knot-ensemble pressure and stress closure}
    \label{subsec:rt_maxwell_knot_stress}

\textbf{[ORTHODOX KINETIC FORM / SST BRIDGE].}
The momentum-flux construction follows the kinetic-theory stress logic used in Maxwellian gas theory \cite{Maxwell1860GasesI}.
Let
\begin{equation}
    \mathbf c
    :=
    \mathbf P-M_K\mathbf U
    \label{eq:rt_maxwell_peculiar_momentum}
\end{equation}
be the peculiar momentum of a knot relative to the ensemble mean flow.  The
kinetic contribution to the material stress is
\begin{equation}
    \Pi_{ij}^{\rm knot,kin}
    =
    \sum_K
    \int
    \frac{c_i c_j}{M_K}
    f_K\,d\Gamma_K,
    \label{eq:rt_maxwell_knot_stress_tensor}
\end{equation}
where $d\Gamma_K$ denotes the complete reduced phase-space measure.  Since
$[c_i c_j/M_K]={\rm J}$ and the integrated distribution has units
${\rm m^{-3}}$, $[\Pi_{ij}]={\rm J\,m^{-3}}={\rm Pa}$.

The isotropic knot-ensemble pressure is
\begin{equation}
    p_{\rm knot}
    =
    \frac13
    \operatorname{tr}
    \Pi^{\rm knot,kin},
    \label{eq:rt_maxwell_knot_pressure}
\end{equation}
and the full material closure may contain additional interaction and internal
mode stresses,
\begin{equation}
    \Pi_{ij}^{\rm knot}
    =
    \Pi_{ij}^{\rm knot,kin}
    +
    \Pi_{ij}^{\rm knot,int}
    +
    \Pi_{ij}^{\rm knot,mode}.
    \label{eq:rt_maxwell_full_knot_stress}
\end{equation}

As a scale check, the canonical SST substrate dynamic-pressure scale is
\begin{align}
    p_{\rm substrate,0}
    &=
    \frac12\rhoF\vchar{}^2\nonumber\\
    &=
    4.1877439\times10^5\,\mathrm{Pa},
    \label{eq:rt_maxwell_substrate_pressure_scale}
\end{align}
using $\rhoF=7.0\times10^{-7}\,\mathrm{kg\,m^{-3}}$ and
$\vchar=1.09384563\times10^6\,\mathrm{m\,s^{-1}}$.  This is a
substrate Euler/Bernoulli scale, not a prediction for $p_{\rm knot}$.

\textbf{[CANON GUARD].}
The knot-ensemble pressure $p_{\rm knot}$ is distinct from the Euler/Bernoulli
pressure of the SST substrate.  They may exchange momentum through a derived
coupling law, but they must not be identified or added twice in the same energy
or stress ledger.

\subsection{Priority IV: isotropization of anisotropic knots}
    \label{subsec:rt_maxwell_orientation_isotropy}

Let $g_K(R,t)$ be the orientation marginal on $SO(3)$, normalized with respect
to the Haar measure $dR$,
\begin{equation}
    \int_{SO(3)} g_K(R,t)\,dR=1.
    \label{eq:rt_maxwell_orientation_normalization}
\end{equation}
For the conventional Euler-angle normalization
$\int_{SO(3)}dR=8\pi^2$, the isotropic distribution is
\begin{equation}
    g_K^{\rm iso}(R)
    =
    \frac{1}{8\pi^2}.
    \label{eq:rt_maxwell_orientation_isotropic}
\end{equation}

A convenient second-rank anisotropy diagnostic is
\begin{equation}
    Q_{ij}^{(K)}
    =
    \left\langle
        t_i t_j-\frac{1}{3}\delta_{ij}
    \right\rangle_{g_K},
    \label{eq:rt_maxwell_orientation_tensor}
\end{equation}
for a declared body-fixed director $\mathbf t$.  Ensemble isotropization requires
\begin{equation}
    Q_{ij}^{(K)}(t)
    \longrightarrow
    0
    \label{eq:rt_maxwell_isotropization_gate}
\end{equation}
under the same interaction kernel used in
Eq.~\eqref{eq:rt_maxwell_sst_boltzmann}.

\textbf{[CRITICAL NOTE].}
Equation~\eqref{eq:rt_maxwell_isotropization_gate} concerns an ensemble of
randomly oriented knots.  It does not prove that the field of a single isolated
anisotropic knot is locally isotropic.  Single-particle isotropy and ensemble
isotropy are separate theorem targets.

\subsection{Numerical programme and preregistered observables}
    \label{subsec:rt_maxwell_numerical_programme}

For each candidate knot $K$, the minimal numerical campaign is:
\begin{enumerate}
    \item obtain a converged relaxed finite-core reference state
    $\mathcal Q_K^{(0)}$;
    \item linearize the declared Euler/Biot--Savart/finite-core model and remove
    rigid translation and global-orientation zero modes before classifying the
    internal spectrum;
    \item compute $\omega_{K,n}$ and classify each eigenvector by Kelvin/shape,
    framed twist, writhe projection, and core-profile content;
    \item determine whether a true $\Delta_{K,n}>0$ exists independently of
    $\omega_{K,n}>0$;
    \item perform controlled two-knot and externally driven encounters and
    measure the projected transfer matrix $G^{KL}_{ab}$;
    \item measure
    \begin{equation}
        \Delta E_{\rm CM},\quad
        \Delta E_{\rm rot/orient},\quad
        \Delta E_{\rm Kelvin},\quad
        \Delta E_{\rm twist},\quad
        \Delta E_{\rm core},
        \label{eq:rt_maxwell_energy_transfer_channels}
    \end{equation}
    while separately recording $\Delta Wr$ as a geometric response rather than
    automatically treating it as an independent energy channel;
    \item infer $\sigma^{KL}_{a,\rm inel}(E,\Omega)$ or a nonlocal replacement
    when a cross-section description fails;
    \item compute $C_{V,K}^{\rm int}(T)$, predicted inelastic thresholds,
    line shifts/broadening and $\Pi_{ij}^{\rm knot}$ without retuning the mode
    spectrum;
    \item test orientation relaxation
    $Q_{ij}^{(K)}(t)\to0$ on held-out encounter ensembles.
\end{enumerate}

\subsection{Primary falsifiers}

\begin{enumerate}
    \item \textbf{Thermodynamic mode-count failure.}
    A strongly coupled, experimentally accessible, equilibrated internal mode is
    predicted but its heat-capacity/inelastic signature is absent.

    \item \textbf{Spectroscopic failure.}
    The calculated occupation and coupling of internal modes imply line shifts,
    sidebands, or widths beyond experimental limits.

    \item \textbf{Gap failure.}
    A claimed positive gap is inferred only from $\omega_{K,n}>0$ while the
    underlying classical energy functional permits $E-E_0\to0$ continuously.

    \item \textbf{Closure failure.}
    The same interaction kernel cannot reproduce both the measured mode-energy
    transfer matrix and the coarse-grained stress/isotropization observables.

    \item \textbf{Resolution failure.}
    The classified low-energy spectrum or coupling hierarchy does not converge
    under refinement of centerline resolution, core resolution, regularization,
    and domain size.
\end{enumerate}

\subsection{Research target statement}

\begin{center}
\begin{minipage}{0.92\linewidth}
\itshape
For every candidate knot $K$, SST must derive a mode-resolved interaction
ledger that distinguishes collective translation and orientation from genuine
internal Kelvin, twist and core excitations, treats writhe as a geometric
projection unless independently dynamical, determines which modes are coupled,
kinematically accessible and equilibrated, and propagates those results without
retuning into thermodynamic, spectroscopic, stress and isotropization
predictions.  A candidate that fails any held-out empirical gate is rejected or
its dynamical closure must be revised.
\end{minipage}
\end{center}
